\documentclass[../main.tex]{subfiles}

\begin{document}

\begin{thm}[Squeeze] \label{thm:limit-squeeze}
  Let {\({\color{main}f(x)} \le g(x) \le {\color{supp}h(x)}\)} when \(x\) is near \(a\) and 
  \[ 
    \lim_{x \to a} {\color{main}f(x)} = L \quad\text{and}\quad \lim_{x \to a} {\color{supp}h(x)} = L,
  \]
  then
  {\(\lim_{x \to a} g(x) = L\).}

  \textbook{Theorem 2.7 on page 171}
\end{thm}


\begin{sidefigure}{%
  \includestandalone[page=6]{../standalones/mindmap-limits}}
\begin{example} \label{ex:limit-law-x-square-sin(x)}
  Evaluate \(\lim_{x \to 0} x^{2} \sin\left(\frac{1}{x}\right)\).

  \blanks*
\end{example}
\end{sidefigure}
\clearpage

This exercise is very similar to Example~\ref{ex:limit-law-x-square-sin(x)}. We walk you through two approaches to evaluate the limit. Both approaches are valid. Pick your favourite.

\begin{exercise} \label{ex:limit-law-x-sin(x)}
  Evaluate \(\lim_{x \to 0} x \sin\left(\frac{1}{x}\right)\).

  \medskip
  \begin{minipage}{2in}
    \includestandalone[width=2in]{../standalones/plot_squeeze}
  \end{minipage}
  \hspace{0.5in}
  \begin{minipage}{4.5in}
    This problem is only a tad different from Example~\ref{ex:limit-law-x-square-sin(x)}. However, its solution is a bit more complicated. There are two approaches:
    \begin{enumerate}[{label=(\Alph*)}, itemsep={1ex}, topsep={1ex}]
      \item Use piecewise functions. 
      \item Evaluate two one-sided limits if you'd rather not deal with piecewise functions.
    \end{enumerate}

    We need the squeeze theorem in both approaches.  
  \end{minipage}

\textbf{Approach (A).}
\begin{enumerate}[wide]
  \item The form of \(x \sin(1/x)\) suggests that we choose
    \begin{itemize}
      \item the lower bound to be \(f(x) = \underline{\hspace{1in}}\),
      \item the upper bound to be \(h(x) = \underline{\hspace{1in}}\).
    \end{itemize}

  \item Sketch your choices of \(f(x)\) and \(h(x)\) on the graph above. Do they \enquote{squeeze} \(x \sin(1/x)\)? If so, move on to the next step; otherwise, go back to the previous step and try again. Hint: \(f(x) = -x\) and \(h(x) = x\) are the wrong choices.

  \item The piecewise nature of \(f(x)\) and \(h(x)\) forces us to show \(f(x) \le x \sin(1/x) \le h(x)\) in two steps. \updatenote{"g(x)" corrected to "h(x)" to match the upper bound named in the steps below.}
    \begin{enumerate}
      \item Show that \(f(x) \le x \sin\left(\frac{1}{x}\right) \le h(x)\) \hlwarn{for \(x \ge 0\)}.

        \blanklines{10}

      \item Show that \(f(x) \le x \sin\left(\frac{1}{x}\right) \le h(x)\) \hlwarn{for \(x < 0\)}.

        \blanklines{10}
    \end{enumerate}

  \item Apply the Squeeze Theorem to finish evaluating the limit.

    \blanklines{4}
\end{enumerate}
\clearpage

\includestandalone[width=2in]{../standalones/plot_squeeze}

\textbf{Approach (B).} To avoid dealing with piecewise functions, we can calculate two one-sided limits.

\begin{enumerate}[wide]
  \item Use the Squeeze Theorem to evaluate \(\lim_{x \to 0^{-}} x \sin(1/x)\).

    \blanklines{20}

  \item Use the Squeeze Theorem to evaluate \(\lim_{x \to 0^{+}} x \sin(1/x)\). \updatenote{Second limit corrected from \(0^{-}\) to \(0^{+}\) so the two one-sided limits are actually different.}

    \blanklines{20}

  \item The two one-sided limits should have the same value. Conclude that \(\lim_{x \to 0} x \sin(1/x) = \underline{\hspace{1cm}}\).
\end{enumerate}
\end{exercise}

The next exercises are direct application of the squeeze theorem. No need to choose an upper and a lower bound for \(g(x)\) because they are already given. \updatenote{"an lower" corrected to "a lower".} It is more of a ``do I remember the statement of the squeeze theorem?'' type of exercise.
\begin{exercise} \label{ex:limit-squeeze-easy}
  Suppose \(x+1 \le g(x) \le x^{3} - 2x + 3\). Find \(\lim_{x \to 1} g(x)\).

  \blanks*
\end{exercise}

This requires a combination of the techniques from Exercise~\ref{ex:limit-law-x-square-sin(x)} and the one above.
\begin{exercise}[Midterm, Fall 2024]
  Evaluate \(\lim_{x \to 0} f(x)\) given that \(0 \le f(x) \le |x \sin(1/x)|\) for all \(x\).

  \blanks*
\end{exercise}
\clearpage
\end{document}
